\section{Methods}
\label{sec:methods}

This section describes the six models (Section~\ref{sec:models}), the
demonstrations and how we split them (Section~\ref{sec:data}), how we measure
each model (Section~\ref{sec:measures}), and how we crop the tactile images
(Section~\ref{sec:cropmethod}).

\subsection{The models}
\label{sec:models}

\begin{newtext}
Every model reads the joint positions of both arms (twelve numbers, which we
call \emph{proprioception}) and some of the cameras. Each outputs a
\emph{chunk}: a plan of future arm commands at \SI{30}{\hertz}.
Table~\ref{tab:models} lists how they differ. For each model we give its
original design first and then what we changed.

\textbf{ACT} \citep{zhao2023act} encodes each camera with a ResNet-18 and
decodes a whole chunk at once with a transformer, trained as a conditional
variational autoencoder. The original smooths overlapping chunks at test time.
We use the LeRobot implementation unchanged, with all five cameras and a chunk
of 100 steps (\SI{3.3}{\second}).

\textbf{Diffusion Policy} \citep{chi2023diffusion} starts from random noise
and removes it step by step to produce a chunk, using a one-dimensional
convolutional network conditioned on the images. We use the LeRobot version
with all five cameras and ResNet-18 encoders trained from scratch. It reads the
last two frames, predicts 64 steps and executes 32, so we score the 32 it
executes.

\textbf{pi0.5} \citep{pi05} is a vision-language-action model: a pretrained
vision-language model reads the images and a written instruction, and a
smaller ``action expert'' produces the chunk by flow matching. It was
pretrained on a large set of robot data. We finetune the public base model with
LoRA \citep{hu2022lora} (rank 16), which trains a small low-rank addition to
the action expert and leaves the rest frozen. It has only three image slots, so
it reads the overhead camera and one fingertip per arm. It is given the
instruction ``fold the short'' and plans 50 steps.

\textbf{Flow matching} is our own small policy. It keeps pi0.5's training
objective and action expert but replaces the pretrained vision-language model
with Diffusion Policy's ResNet-18 encoders, trained from scratch on all five
cameras. Flow matching \citep{lipman2023flow} learns to move noise to an action
chunk along straight lines, which is the optimal-transport path between the
two. The model therefore separates the objective from pi0.5's pretraining.

\textbf{DreamZero} \citep{dreamzero} is a \emph{world action model}: it
predicts future video and future actions together, under one flow-matching
objective. The original is built on a 14-billion-parameter pretrained video
generator. Ours has about 59 million parameters and is trained from scratch on
our demonstrations. It tiles all five cameras into one image and compresses
each frame with a frozen image autoencoder. It sees two past frames, up to
\SI{1.6}{\second} back, and predicts frames every \SI{0.8}{\second} up to
\SI{4.8}{\second} ahead, with 24 action steps per chunk.

\textbf{FastWAM} is also a world action model, built on a pretrained
five-billion-parameter video generator (Wan2.2) with an action expert beside
it. We finetune the public LeRobot base model. Its image input is one frame of
fixed size, so it reads the overhead camera and a two-by-two composite of the
four fingertip images. It predicts nine frames over about a second and plans 32
action steps, of which it executes 10.
\end{newtext}

\begin{table}[t]
  \centering
  \caption{\new{The six models. \emph{Plan} is the number of steps each model
  commits to and we score (at \SI{30}{\hertz}). \emph{Fingertips} is how many
  of the four tactile cameras the model reads. Parameters count the whole
  model; for pi0.5 only the LoRA addition is trained.}}
  \label{tab:models}
  \small
  \begin{tabular}{lrrlrrl}
    \toprule
    model & plan & fingertips & starts from & parameters & steps $\times$ batch & predicts images \\
    \midrule
    ACT & 100 & 4 & ImageNet encoder & \SI{52}{M} & \num{80000} $\times$ 8 & no \\
    Diffusion & 32 & 4 & scratch & \SI{323}{M} & \num{100000} $\times$ 8 & no \\
    pi0.5 & 50 & 2 & robot-pretrained & \SI{4.1}{B} & \num{30000} $\times$ 1 & no \\
    Flow matching & 50 & 4 & scratch & \pending{count} & \num{60000} $\times$ 16 & no \\
    DreamZero & 24 & 4 & scratch & \SI{59}{M} & \num{80000} $\times$ 8 & yes \\
    FastWAM & 10 & 4 (composite) & video-pretrained & \SI{6.0}{B} & \num{30000} $\times$ 8 & yes \\
    \bottomrule
  \end{tabular}
\end{table}

\subsection{The demonstrations and the split}
\label{sec:data}

We recorded one task in one session: a pair of shorts, already laid flat, is
folded once by the two arms. There are 65 demonstrations, about 15 seconds
each, 30141 frames in all at \SI{30}{\hertz}. \new{Figure~\ref{fig:trajectory}
shows one demonstration from start to end.}

\begin{figure}[t]
  \centering
  \includegraphics[width=0.78\textwidth]{trajectory_strip.png}
  \caption{\new{One held-out demonstration seen by the overhead camera, 20
  frames evenly spaced from start to end. The arms grasp the waistband, fold
  the shorts over and release.}}
  \label{fig:trajectory}
\end{figure}

We train every model on the first 58 demonstrations and test it on the
\emph{last seven}, which it never sees. \new{The last seven were recorded at
the end of the session. If something changed slowly during the session, such
as the lighting or wear on the gels, the test set would differ from the
training set for that reason alone. To check this, we retrain ACT and
Diffusion on a \emph{random seven} held out from across the whole session.
Section~\ref{sec:compare} compares the two.}

\subsection{How we measure}
\label{sec:measures}

\textbf{Action error.} We give a model one observation and compare the chunk it
plans with the commands the demonstrator actually sent next. We take the
root-mean-square error (RMSE) over all twelve joints, in degrees, separately for
each step of the chunk. We sample one frame in twenty, which gives 1127 scored
frames from the training demonstrations and 111 from the test demonstrations.
Models plan different distances ahead, and error grows with distance, so we
compare models over the first ten steps (\SI{0.33}{\second}), which every model
plans. The \emph{gap} is the test error divided by the training error. Models
that sample have their random seed fixed, so a score repeats exactly.

\begin{newtext}
\textbf{Grad-CAM} \citep{selvaraju2017gradcam} shows where in an image a
model's output comes from. We take the magnitude of the planned chunk as the
output. We then compute its gradient with respect to the last feature map of
the ResNet encoder, average that gradient over each channel to get one weight
per channel, and sum the channels with these weights. Negative values are set
to zero, and the map is scaled so its peak is one and enlarged to the image
size. Grad-CAM needs a convolutional feature map, so it applies to ACT and
Diffusion but not to pi0.5.

\textbf{Input contribution.} To ask how much a model uses one input, we replace
that input with an uninformative version and measure how far the plan moves.
An image is replaced by a flat image of its own average colour, which keeps its
brightness but removes all structure. The joint positions are replaced by a
constant. The distance between the two plans is the mean, over the chunk, of the
Euclidean distance between their commands. We do this for each input in turn
and divide each distance by their sum, so the contributions of one model add up
to 100\,\%. Because they are normalised within a model, the percentages rank a
model's inputs; they cannot be compared in size between models. For models
that sample, the seed is fixed during this, since otherwise sampling noise
alone moves the plan by far more than removing a camera does.

\textbf{World models.} We score frames in two ways. \emph{Reconstruction}
passes frames the model has already seen through its own image autoencoder and
back. \emph{Prediction} compares the frames the model generates for the future,
together with its own planned actions, with what the cameras actually recorded.
Both use peak signal-to-noise ratio (PSNR, in decibels, higher is better) and
the structural similarity index (SSIM, from 0 to 1). We score every sensor
separately. As a baseline we repeat the last frame the model saw; a model that
cannot beat this has learned nothing about how the scene moves.
\end{newtext}

\subsection{The tactile crop}
\label{sec:cropmethod}

\begin{newtext}
Our earlier Grad-CAM pictures showed ACT putting weight near the edges of the
tactile images before contact. Figure~\ref{fig:rimevidence} measures what the
images themselves show before contact. On each sensor the brightest edge is 10 to
21\,\% brighter than the centre, which fits light leaking in at the gel boundary. However, the
brightest feature on every sensor is a vertical ridge inside the gel, about a
third of the way across, not the edge. So the images support a brighter rim,
but they do not show that the rim is the strongest signal before contact.
\end{newtext}

\begin{figure}[t]
  \centering
  \includegraphics[width=0.92\textwidth]{rim_evidence.png}
  \caption{\new{The four fingertip sensors before contact. Top: the first frame
  of one demonstration, with the rows-only crop (dashed) and the four-edge crop
  (solid). Bottom: brightness of each row and column relative to the centre of
  the gel, averaged over the first frame of all 65 demonstrations; dotted lines
  mark the crop boundaries. The sharp peak in the columns is a ridge in the gel,
  not the rim.}}
  \label{fig:rimevidence}
\end{figure}

We test three crops, each applied to all four fingertip images. \emph{Rows
only} removes the top and bottom tenth of the image. We measured beforehand
that the columns carry contact signal right up to the edge on two sensors,
while every sensor has spare rows. \new{\emph{Four edges} removes a tenth from
every side.} Both then resize the image back to its original shape, which
stretches it. \emph{Rows without stretching} keeps the cropped image at its
smaller size. Only ACT accepts cameras of different sizes, so only ACT gets
this third crop. \new{The crop is part of the model's saved input pipeline, so
it is applied during training and again at test time. We confirmed this in the
saved pipeline of every cropped model.} Apart from the crop, a cropped model is
identical to its uncropped version: the same architecture, data, number of
steps, batch size and random seed.
